\begin{frame}[fragile]{Using UHD from Python (2/7)}
\vspace{-0.5cm}
\footnotesize

\begin{itemize}
    \item \textbf{\textcolor{DarkBlue}{Step 1: Enable the UHD Python API}}
\end{itemize}

\hspace{2.8em}
\parbox{0.88\linewidth}{
\footnotesize
While configuring UHD, enable the Python API by passing
\texttt{-DENABLE\_PYTHON\_API=ON} to CMake. This builds the UHD Python
module along with the UHD C++ library.
}

\vspace{0.2cm}

\scriptsize
\begin{Verbatim}[xleftmargin=3.8em,breaklines,breaksymbolleft={},breaksymbolright={}]
cmake ../ -DCMAKE_BUILD_TYPE=Release \
-DENABLE_PYTHON_API=ON \
-DPYTHON_EXECUTABLE=$(which python3)
\end{Verbatim}





\hspace{2.8em}
\parbox{0.88\linewidth}{
\scriptsize
\textbf{\textcolor{DarkBlue}{Verify the Enabled Components}}

During CMake configuration, verify that
\textbf{LibUHD - Python API} appears under the enabled components.
This confirms that the UHD Python bindings will be built.
}
\vspace{0.5cm}
\begin{itemize}
    \item \textbf{\textcolor{DarkBlue}{Step 2: Build and Install UHD}}
\end{itemize}

\hspace{2.8em}
\parbox{0.88\linewidth}{
\footnotesize
After enabling the Python API, compile UHD and install it system-wide.
This builds the Python bindings along with the UHD libraries.
}

\vspace{0.2cm}

\scriptsize
\begin{Verbatim}[xleftmargin=3.8em]
make -j$(nproc)
sudo make install
\end{Verbatim}

\vspace{0.2cm}



\end{frame}